\documentclass{article}
\usepackage[left=3cm,top=3cm,right=2cm,bottom=2cm]{geometry}
\usepackage{amssymb}
\usepackage{amsmath}
\usepackage{cancel}
\usepackage[hidelinks]{hyperref}

\title{(Universidade de São Paulo)\\
Equações Elípticas}

\author{Métodos Numéricos em Equações Diferenciais\\
Docente: Fabrício Simeoni de Sousa\\
Vinícius de Sá Ferreira -- 15491650}

\date{\today}

\begin{document}
    \maketitle
    
    \section{Problema}
    \begin{align}
        &\frac{\partial^2 u}{\partial x^2}(x,y) + 3\frac{\partial^2 u}{\partial y^2}(x,y) = (4\pi^2-3)\frac{\partial u}{\partial y}(x,y),& (x,y) \in [0,1] \times [-1,1] \label{eq}\\
        &u(0,y)=u(1,y)=0,& y \in [-1,1] \label{dirichlet}\\
        &\frac{\partial u}{\partial \mathbf{n}}(x,-1) = e \sin(2\pi x),\ \frac{\partial u }{\partial \mathbf{n}}(x,1) + u(x,1) = 0,& x \in [0,1] \label{neumman}
    \end{align}

    Em \ref{neumman} temos
    \[\mathbf{n}(x,y) = \begin{pmatrix}0 \\ y\end{pmatrix},\ (x,y) \in [0,1] \times \{-1,1\}\]
    logo
    \[\frac{\partial u}{\partial \mathbf{n}}(x,-1) = e \sin(2\pi x) \iff
    \begin{pmatrix}\frac{\partial^2 u}{\partial x^2}(x,-1) & \frac{\partial^2 u}{\partial y^2}(x,-1)\end{pmatrix} \begin{pmatrix}0 \\ -1\end{pmatrix} = e \sin(2\pi x)\]
    \[\iff -\frac{\partial^2 u}{\partial y^2}(x,-1) = e \sin(2\pi x)\]
    e
    \[\frac{\partial u }{\partial \mathbf{n}}(x,1) + u(x,1) = 0 \iff
    \begin{pmatrix}\frac{\partial^2 u}{\partial x^2}(x,1) & \frac{\partial^2 u}{\partial y^2}(x,1)\end{pmatrix} \begin{pmatrix}0 \\ 1\end{pmatrix} + u(x,1) = 0\]
    \[\iff \frac{\partial^2 u}{\partial y^2}(x,1) + u(x,1) = 0\]

    Portanto, \ref{neumman} pode ser reescrito como
    \begin{align}
        &-\frac{\partial^2 u}{\partial y^2}(x,-1) = e \sin(2\pi x),\ \frac{\partial^2 u}{\partial y^2}(x,1) + u(x,1) = 0,& x \in [0,1] \label{neumman_remake}
    \end{align}

    \section{Discretização}
    Usamos os métodos de diferenças finitas centrados para primeira e segunda derivadas, que são $\mathcal{O}(h^2)$
    \[u'(x_n) \approx Du(x_n) = \frac{u(x_{n+1})-u(x_{n-1})}{2h}\]
    \[u''(x_n) \approx D^2u(x_n) = \frac{u(x_{n-1})-2u(x_n)+u(x_{n+1})}{h^2}\]
    Definimos $x_i = ih$ e $y_j = -1 + jh$, com $i \in \{0,1,\cdots,n+1\}$ e $j \in \{0,1,\cdots,m+1\}$, onde $h = \frac{1}{n+1} = \frac{2}{m+1}$. Então o problema se transforma em
        \begin{align}
        &\frac{U_{i-1,j}-2U_{i,j}+U_{i+1,j}}{h^2} + 3\left(\frac{U_{i,j-1}-2U_{i,j}+U_{i,j+1}}{h^2}\right) = (4\pi^2-3)\frac{U_{i,j+1}-U_{i,j-1}}{2h} \label{num_eq}\\
        &u(0,y_j) = u(1,y_j) = 0 \label{num_dirichlet}\\
        &-3\left(\frac{U_{i,-1}-2U_{i,0}+U_{i,1}}{h^2}\right) = e \sin(2\pi x),\ \frac{U_{i,m}-2U_{i,m+1}+U_{i,m+2}}{h^2} + U_{i,m+1} = 0 \label{num_neumman}
    \end{align}
\end{document}